\clearpage
\VerdierPhysicalPageAnchor{0122}%
\begingroup
\small
\setlength{\emergencystretch}{1.5em}%
\setlength{\abovedisplayskip}{4pt plus 2pt minus 2pt}%
\setlength{\belowdisplayskip}{4pt plus 2pt minus 2pt}%

\VerdierSemanticAnchor{ch2-1-2-13-countability-hypothesis}
\textit{如果可数积或可数直和在 $\mathcal D$ 中可表示，则范畴
$\mathcal D$ 可分解。}

\VerdierSemanticAnchor{ch2-1-2-13-proof-morphism}
事实上，设 $v\colon W\longrightarrow X$ 满足 $fv=0$。于是有态射范畴中的
一个态射：

\VerdierSemanticAnchor{ch2-1-2-13-morphism-square}
\[
  \begin{tikzcd}[column sep=4em,row sep=2.5em]
    X \arrow[r,"f"] & Y \\
    W \arrow[u,"v"] \arrow[r] & 0 \arrow[u]
  \end{tikzcd}
  \,.
\]

\VerdierSemanticAnchor{ch2-1-2-13-1}
因此，应用函子 $\operatorname{Tr}$，得到交换图：
\begin{equation}
  \resizebox{0.82\linewidth}{!}{%
    $\begin{tikzcd}[ampersand replacement=\&,column sep=2.05em,row sep=3em]
      X \arrow[r,"f"]
        \& Y \arrow[r,"{a(f)}"]
        \& \tau(f) \arrow[r,"{b(f)}"]
        \& X[1] \\
      W \arrow[u,"v"] \arrow[r]
        \& 0 \arrow[u] \arrow[r]
        \& \tau(W\longrightarrow 0) \arrow[u,"{\operatorname{Tr}(v,0)}"]
          \arrow[r,"{b(W\longrightarrow 0)}"']
        \& W[1] \arrow[u,"{v[1]}"']
    \end{tikzcd}$%
  }\,.
  \tag{1.2.13.1}
\end{equation}

\VerdierSemanticAnchor{ch2-1-2-13-bottom-triangle}
先注意图 (1.2.13.1) 的下行是特异三角。因此，态射
$b(W\longrightarrow 0)$ 是同构 (1.2.1)。令：

\VerdierSemanticAnchor{ch2-1-2-13-alpha-definition}
\[
  \alpha(v)=\operatorname{Tr}(v,0)[-1]
  b(W\longrightarrow 0)^{-1}[-1]\,,
\]
则下图交换：

\VerdierSemanticAnchor{ch2-1-2-13-alpha-diagram}
\[
  \begin{tikzcd}[column sep=3.5em,row sep=2.4em]
    \tau(f)[-1] \arrow[r,"{b(f)[-1]}"]
      & X \arrow[r,"f"]
      & Y \\
      & W \arrow[ul,"{\alpha(v)}"'] \arrow[u,"v"] \arrow[ur,"0"']
  \end{tikzcd}
\]
立即验证可知，这样定义了一个截面：
\par

\VerdierSemanticAnchor{ch2-1-2-13-alpha-section}
\[
  \alpha\colon \operatorname{Ker}(f)\longrightarrow \tau(f)[-1]
\]
\noindent 即由 $\tau(f)[-1]$ 所表示的函子到态射 $f$ 的核函子的典范态射
$\tau(f)[-1]\longrightarrow\operatorname{Ker}(f)$ 的一个截面。

\VerdierSemanticAnchor{ch2-1-2-13-projector-kernel}
因此，$f$ 的核函子同构于 $\tau(f)[-1]$ 的某个投影态射的核函子，
故可表示 (1.2.9)。

\VerdierSemanticAnchor{ch2-1-2-13-proof-complete}
由 (1.2.9) 可知态射 $f$ 可分解，证明完毕。
\endgroup
